% Figure: interpolating a tabulated device curve. A high-degree global
% polynomial through the nodes oscillates between them (Runge); a cubic
% spline through the same nodes stays close to the underlying curve.
% Two panels side by side share the node set.
\begin{figure}[htbp]
\centering
\resizebox{0.98\linewidth}{!}{%
\begin{tikzpicture}[scale=1.0,
  ax/.style={->, draw=engblack, thick},
  node/.style={circle, fill=engblack, inner sep=1.2pt},
  truec/.style={draw=engblack!45, thick, dash pattern=on 2pt off 2pt},
  badc/.style={draw=engred, very thick},
  goodc/.style={draw=engblue, very thick},
  ann/.style={font=\footnotesize, align=center},
]
  % The underlying device curve is f(u) = 1/(1 + 12 u^2) on u in [-1,1],
  % the Runge function; nodes equally spaced. Display: x = 2.4*(u+1)/2,
  % y = 2.6*f. Two panels offset horizontally.
  \def\sx{2.4}\def\sy{2.6}
  % ---- Panel (a): global polynomial ----
  \begin{scope}[shift={(0,0)}]
    \draw[ax] (-0.2, 0) -- ({\sx + 0.4}, 0) node[below right, ann] {input};
    \draw[ax] (0, -1.0) -- (0, {\sy + 0.5}) node[left, ann] {output};
    % True curve.
    \draw[truec, domain=-1:1, samples=80, smooth]
      plot ({\sx*(\x+1)/2}, {\sy*(1/(1 + 12*\x*\x))});
    % Oscillating interpolant (schematic high-degree polynomial):
    % a curve that hits the nodes but swings negative near the ends.
    \draw[badc, smooth, tension=0.9]
      plot coordinates {
        ({\sx*0.00}, {\sy*0.077})
        ({\sx*0.07}, {-\sy*0.20})
        ({\sx*0.125}, {\sy*0.16})
        ({\sx*0.25}, {\sy*0.25})
        ({\sx*0.375}, {\sy*0.32})
        ({\sx*0.50}, {\sy*1.0})
        ({\sx*0.625}, {\sy*0.32})
        ({\sx*0.75}, {\sy*0.25})
        ({\sx*0.875}, {\sy*0.16})
        ({\sx*0.93}, {-\sy*0.20})
        ({\sx*1.00}, {\sy*0.077}) };
    % Nodes (6 equally spaced points on the true curve).
    \foreach \u in {-1, -0.6, -0.2, 0.2, 0.6, 1} {
      \node[node] at ({\sx*(\u+1)/2}, {\sy*(1/(1 + 12*\u*\u))}) {};
    }
    \node[ann, anchor=north] at ({\sx*0.5}, -0.55)
      {(a) global polynomial:\\swings between nodes};
  \end{scope}
  % ---- Panel (b): cubic spline ----
  \begin{scope}[shift={(4.4,0)}]
    \draw[ax] (-0.2, 0) -- ({\sx + 0.4}, 0) node[below right, ann] {input};
    \draw[ax] (0, -1.0) -- (0, {\sy + 0.5}) node[left, ann] {output};
    \draw[truec, domain=-1:1, samples=80, smooth]
      plot ({\sx*(\x+1)/2}, {\sy*(1/(1 + 12*\x*\x))});
    % Cubic spline: smooth curve hugging the true curve through the nodes.
    \draw[goodc, smooth, tension=0.7]
      plot coordinates {
        ({\sx*0.00}, {\sy*0.077})
        ({\sx*0.20}, {\sy*0.18})
        ({\sx*0.40}, {\sy*0.45})
        ({\sx*0.50}, {\sy*1.0})
        ({\sx*0.60}, {\sy*0.45})
        ({\sx*0.80}, {\sy*0.18})
        ({\sx*1.00}, {\sy*0.077}) };
    \foreach \u in {-1, -0.6, -0.2, 0.2, 0.6, 1} {
      \node[node] at ({\sx*(\u+1)/2}, {\sy*(1/(1 + 12*\u*\u))}) {};
    }
    \node[ann, anchor=north] at ({\sx*0.5}, -0.55)
      {(b) cubic spline:\\stays near the curve};
  \end{scope}
\end{tikzpicture}%
}%
\caption[Spline against global polynomial for interpolating a device
curve]{Two interpolants through the same six tabulated points (black)
on a peaked device curve (grey dashed). Panel (a): a single
high-degree polynomial forced through every node hits the nodes but
oscillates between them, swinging below zero near the ends, the Runge
phenomenon. Panel (b): a cubic spline, a piecewise-cubic function
joined so that value, slope, and curvature match at each node, stays
close to the underlying curve everywhere. For interpolating a measured
device characteristic the spline is the working tool, because the
engineer cares about the value between the nodes and a global
polynomial buys exact node values at the cost of nonsense between
them.}
\label{fig:vol02:ch02:spline-vs-polynomial}
\end{figure}
